\documentclass[../../main.tex]{subfiles}% 注意这里的文件路径不能用 ./main.tex ,否则用latexmk编译子文件会报错
\graphicspath{{\subfix{./image/}}{\subfix{../../image/}}} % 指定图片引用路径,后续可以直接使用图片文件名
% 如果图片存放在上级目录,则使用 ../../image/ 作为备用路径

% 例如:
% \begin{figure}[H]
% \centering
% \includegraphics[width=0.3\textwidth]{图.png}
% \caption{}
% \label{figure:图}
% \end{figure}
% 注意:上述\label{}一定要放在\caption{}之后,否则引用图片序号会只会显示??.

\begin{document}

\section{极限 Young 测度的刻画}\label{section:ch08-测度刻画}

设 $(u,\nu)$ 为\cref{theorem:ch08-存在性} 给出的问题 (P) 的整体 Young 测度值解。本节给出 $N=1$ 时 Young 测度 $\nu$ 的显式表达式（与此相关，见\cref{section:ch05-收敛}）。

先陈述若干预备事实。设 $\phi$ 改变单调特征的次数有限；由假设\eqref{align:ch08-8.1c}--\eqref{align:ch08-8.1d} 注意该次数为偶数，记为 $2n$（$n\in\mathbb{N}$）。故
\begin{align}
\mathbb{R}=\left(\bigcup_{l=0}^n I_l\right)\cup\left(\bigcup_{m=1}^n\widehat I_m\right)\quad(n\in\mathbb{N}), \label{align:ch08-8.96}
\end{align}
其中
\begin{align*}
I_0\triangleq (-\infty,b_1],\quad I_l\triangleq (a_l,b_{l+1}]\ (l=1,\cdots ,n-1),\quad I_n\triangleq (a_n,\infty),
\end{align*}
\begin{align*}
\widehat I_m\triangleq (b_m,a_m]\ (m=1,\cdots ,n),
\end{align*}
$b_l$ 为 $\phi$ 图像的局部极大点，$a_l$ 为局部极小点。令 $J_l\triangleq \phi(I_l)$、$\widehat J_m\triangleq \phi(\widehat I_m)$（$l=0,\cdots ,n$；$m=1,\cdots ,n$）；由于 $\phi$ 在每个区间 $I_l$ 上递增、在每个区间 $\widehat I_m$ 上递减，有
\begin{align*}
J_0\triangleq (-\infty,\phi(b_1)],\quad J_l\triangleq (\phi(a_l),\phi(b_{l+1})]\ (l=1,\cdots ,n-1),\quad J_n\triangleq (\phi(a_n),\infty),
\end{align*}
\begin{align*}
\widehat J_m\triangleq [\phi(a_m),\phi(b_m))\ (m=1,\cdots ,n).
\end{align*}
另考虑 $n+1$ 个递增函数 $s_l:J_l\to I_l$、$s_l\triangleq (\phi|_{I_l})^{-1}$（$l=0,\cdots ,n$）与 $n$ 个递减函数 $t_m:\widehat J_m\to\widehat I_m$、$t_m\triangleq (\phi|_{\widehat I_m})^{-1}$（$m=1,\cdots ,n$）。仿照文献~\cite{ref83}，我们用如下假设：
\begin{align*}
\text{函数 }s_0',\cdots ,s_n',t_1',\cdots ,t_n'\text{ 在 }\mathbb{R}\text{ 的任一开子集上线性无关。}
\qquad(C)
\end{align*}

现陈述如下结果。

\begin{theorem}[极限 Young 测度的显式刻画]\label{theorem:ch08-测度刻画}
设 $N=1$，$\Omega\equiv(a,b)$，\eqref{align:ch08-8.1a}--\eqref{align:ch08-8.1d} 成立，$\phi$ 改变单调特征的次数有限。设 $(u,\nu)$ 为\cref{theorem:ch08-存在性} 给出的 (P) 的整体 Young 测度值解。则对几乎所有 $(x,t)\in Q_\infty$，
\begin{align}
\nu_{(x,t)}=\sum_{l=0}^n c_l(x,t)\,\delta_{s_l(w(x,t))}
+\sum_{m=1}^n d_m(x,t)\,\delta_{t_m(w(x,t))}, \label{align:ch08-8.97}
\end{align}
其中 $w\in L^2(0,\infty;H^1(\Omega))\cap L^\infty(Q_\infty)$ 为\eqref{align:ch08-8.45} 中的极限函数。进而 $c_l,d_m\in L^\infty(Q_\infty)$，且对几乎所有 $(x,t)\in Q_\infty$：
\begin{enumerate}[(i)]
\item\label{item:ch08-char-i} $0\leqslant c_l(x,t)\leqslant1$，$0\leqslant d_m(x,t)\leqslant1$；
\item\label{item:ch08-char-ii} $\sum_{l=0}^n c_l(x,t)+\sum_{m=1}^n d_m(x,t)=1$；
\item\label{item:ch08-char-iii} 若 $w(x,t)\leqslant A\triangleq \min\{\phi(a_1),\cdots ,\phi(a_n)\}$，则 $c_0(x,t)=1$；若 $w(x,t)\geqslant B\triangleq \max\{\phi(b_1),\cdots ,\phi(b_n)\}$，则 $c_n(x,t)=1$。
\end{enumerate}
\end{theorem}

\begin{remark}\label{remark:ch08-导数公式}
在\cref{theorem:ch08-测度刻画} 的假设下，由等式\eqref{align:ch08-8.97} 与\cref{remark:ch08-w等于phi星}，对几乎所有 $(x,t)\in Q_\infty$ 有
\begin{align*}
\partial_x u(x,t)=\sum_{l=0}^n c_l(x,t)\,s_l(\phi^*(x,t))+\sum_{m=1}^n d_m(x,t)\,t_m(\phi^*(x,t)).
\end{align*}
\end{remark}

为论证\cref{theorem:ch08-测度刻画}，仿照\cref{theorem:ch6-正则化} 之证明思路。设 $(u,\nu)$ 为\cref{theorem:ch08-存在性} 给出的整体 Young 测度值解，$\{u_{\epsilon_k}\}$ 为构造它所用的序列。由于族 $\{\phi(v_\epsilon)\}$（$v_\epsilon\triangleq \partial_x u_\epsilon$）在 $L^\infty(Q_\infty)$ 中有界（见\eqref{align:ch08-8.26}），由\cref{proposition:ch05-弱收敛} 与对角论证，存在序列 $\{\phi(v_{\epsilon_k})\}$ 与 Young 测度 $\tau\in\mathcal{Y}^+(Q_\infty;\mathbb{R})$，使对所有 $T\in(0,\infty)$ 与 $f\in C_c(\mathbb{R})$，
\begin{align}
f(\phi(v_{\epsilon_k}))\rightharpoonup f^*\ \text{在 }L^1(Q_T),\qquad
f^*(x,t)\triangleq \int_{\mathbb{R}}f(y)\,\mathrm{d}\tau_{(x,t)}(y). \label{align:ch08-8.98}
\end{align}
另一方面，对每个 $f\in C_c(\mathbb{R})$，序列 $\{f(\phi(v_{\epsilon_k}))\}$ 在 $L^1(Q)$ 中有界且一致可积，因而（可能抽子列，不重新标号）由\cref{proposition:ch05-弱收敛} 得
\begin{align}
f(\phi(v_{\epsilon_k}))\rightharpoonup(f\circ\phi)^*\ \text{在 }L^1(Q_T),\qquad
(f\circ\phi)^*(x,t)\triangleq \int_{\mathbb{R}}f(\phi(y))\,\mathrm{d}\nu_{(x,t)}(y). \label{align:ch08-8.99}
\end{align}
由等式\eqref{align:ch08-8.98}--\eqref{align:ch08-8.99} 有
\begin{align}
\tau_{(x,t)}=(\nu_{(x,t)})_\phi\quad\text{对几乎所有 }(x,t)\in Q_\infty, \label{align:ch08-8.100}
\end{align}
即 $\tau_{(x,t)}$ 为 $(\nu_{(x,t)})$ 在 $\phi$ 下的像测度。

\begin{theorem}[$\tau$ 为 Dirac 测度]\label{theorem:ch08-τ为Dirac}
设\cref{theorem:ch08-测度刻画} 的假设满足。设 $\tau\in\mathcal{Y}^+(Q_\infty;\mathbb{R})$ 为\eqref{align:ch08-8.98} 中的 Young 测度。则
\begin{align}
\tau_{(x,t)}=\delta_{w(x,t)}\quad\text{对几乎所有 }(x,t)\in Q_\infty. \label{align:ch08-8.101}
\end{align}
\end{theorem}

对几乎所有 $(x,t)\in Q_\infty$ 与所有 $l=0,\cdots ,n$、$m=1,\cdots ,n$，令
\begin{align}
\sigma^l_{(x,t)}\triangleq (\nu_{(x,t)})_{\phi|_{I_l}},\quad
\widehat\sigma^m_{(x,t)}\triangleq (\nu_{(x,t)})_{\phi|_{\widehat I_m}}, \label{align:ch08-8.102a}\\
\sigma_{(x,t)}\triangleq \sum_{l=0}^n\sigma^l_{(x,t)}+\sum_{m=1}^n\widehat\sigma^m_{(x,t)}. \label{align:ch08-8.102b}
\end{align}
显然，对任意如上的 $(x,t)$ 与 $l,m$，有
\begin{align}
\operatorname{supp}\sigma^l_{(x,t)}\subseteq J_l,\quad
\operatorname{supp}\widehat\sigma^m_{(x,t)}\subseteq\widehat J_m,\quad
\operatorname{supp}\sigma_{(x,t)}\subseteq\phi(\mathbb{R})=\mathbb{R}. \label{align:ch08-8.103}
\end{align}
由\eqref{align:ch08-8.102a}，对每个 $f\in C(\mathbb{R})$ 有
\begin{align}
\int_{J_l}f(y)\,\mathrm{d}\sigma^l_{(x,t)}(y)=\int_{I_l}(f\circ\phi)(y)\,\mathrm{d}\nu_{(x,t)}(y), \label{align:ch08-8.104a}\\
\int_{\widehat J_m}f(y)\,\mathrm{d}\widehat\sigma^m_{(x,t)}(y)=\int_{\widehat I_m}(f\circ\phi)(y)\,\mathrm{d}\nu_{(x,t)}(y). \label{align:ch08-8.104b}
\end{align}
由\eqref{align:ch08-8.104a} 与\eqref{align:ch08-8.102b} 显然 $\sigma_{(x,t)}\in\mathcal{P}(\mathbb{R})$。

由等式\eqref{align:ch08-8.100}，\cref{theorem:ch08-τ为Dirac} 是如下命题的直接推论。

\begin{proposition}[$\sigma$ 为 Dirac 测度]\label{proposition:ch08-σ为Dirac}
设\cref{theorem:ch08-测度刻画} 的假设满足。则对几乎所有 $(x,t)\in Q_\infty$，
\begin{align}
\sigma_{(x,t)}=(\nu_{(x,t)})_\phi,\qquad \sigma_{(x,t)}=\delta_{w(x,t)}. \label{align:ch08-8.105}
\end{align}
\end{proposition}

\begin{remark}\label{remark:ch08-σ第一等式}
\eqref{align:ch08-8.105} 中第一个等式易证。事实上，由\eqref{align:ch08-8.96} 与\eqref{align:ch08-8.104a}，对所有 $f\in C(\mathbb{R})$ 与几乎所有 $(x,t)\in Q_\infty$，
\begin{align*}
\int_{\mathbb{R}}(f\circ\phi)(y)\,\mathrm{d}\nu_{(x,t)}(y)
&=\sum_{l=0}^n\int_{I_l}(f\circ\phi)(y)\,\mathrm{d}\nu_{(x,t)}(y)+\sum_{m=1}^n\int_{\widehat I_m}(f\circ\phi)(y)\,\mathrm{d}\nu_{(x,t)}(y)\\
&=\sum_{l=0}^n\int_{J_l}f(y)\,\mathrm{d}\sigma^l_{(x,t)}(y)+\sum_{m=1}^n\int_{\mathbb{R}}f(y)\,\mathrm{d}\widehat\sigma^m_{(x,t)}(y)
=\int_{\phi(\mathbb{R})}f(y)\,\mathrm{d}\sigma_{(x,t)}(y),
\end{align*}
结论即得。
\end{remark}

借助\cref{theorem:ch08-τ为Dirac}，可证\cref{theorem:ch08-测度刻画}。

\begin{proof}[\cref{theorem:ch08-测度刻画} 的证明]
设 $\{\epsilon_k\}$ 为使\cref{proposition:ch08-收敛N1} 与\eqref{align:ch08-8.45} 的收敛成立的指标序列。对 $k\in\mathbb{N}$，令
\begin{align*}
Q_{k,l}\triangleq \{(x,t)\in Q_\infty\mid v_{\epsilon_k}(x,t)\in I_l\}\quad(l=0,\cdots ,n),
\end{align*}
\begin{align*}
Q_{k,m}\triangleq \{(x,t)\in Q_\infty\mid v_{\epsilon_k}(x,t)\in\widehat I_m\}\quad(m=1,\cdots ,n).
\end{align*}
显然，在 $Q_\infty$ 中几乎处处对所有 $f\in C(\mathbb{R})$ 有
\begin{align}
f(v_{\epsilon_k})=\sum_{l=0}^n\chi_{Q_{k,l}}f((s_l\circ\phi)(v_{\epsilon_k}))
+\sum_{m=1}^n\chi_{Q_{k,m}}f((t_m\circ\phi)(v_{\epsilon_k})). \label{align:ch08-8.106}
\end{align}
对任意固定的 $l,m\in\mathbb{N}$，序列 $\{\chi_{Q_{k,l}}\}$ 与 $\{\chi_{Q_{k,m}}\}$ 在 $L^\infty(Q_\infty)$ 中有界。于是存在二者的子列（为简便不重新标号）与 $c_l,d_m\in L^\infty(Q_\infty)$，使对所有 $T\in(0,\infty)$，
\begin{align}
\chi_{Q_{k,l}}\stackrel{*}{\rightharpoonup}c_l,\quad
\chi_{Q_{k,m}}\stackrel{*}{\rightharpoonup}d_m\ \text{在 }L^\infty(Q_T). \label{align:ch08-8.107}
\end{align}
易见函数 $c_l,d_m$（$l=0,\cdots ,n$；$m=1,\cdots ,n$）具有所述性质。

我们证明，对所有 $l=0,\cdots ,n$ 与 $m=1,\cdots ,n$，在 $L^\infty(Q_T)$（$T\in(0,\infty)$）中有
\begin{align}
\chi_{Q_{k,l}}f((s_l\circ\phi)(v_{\epsilon_k}))\stackrel{*}{\rightharpoonup}c_l\,f(s_l(w)), \label{align:ch08-8.108a}\\
\chi_{Q_{k,m}}f((t_m\circ\phi)(v_{\epsilon_k}))\stackrel{*}{\rightharpoonup}d_m\,f(t_m(w)). \label{align:ch08-8.108b}
\end{align}
另一方面，由于由\eqref{align:ch08-8.25} 族 $\{v_\epsilon\}$ 在 $L^\infty(Q_\infty)$ 中有界因而有界且一致可积，由\cref{proposition:ch05-弱收敛} 对所有 $T\in(0,\infty)$ 有 $f(v_{\epsilon_k})\rightharpoonup f^*\triangleq \int_{\mathbb{R}}f(y)\,\mathrm{d}\nu(y)$ 在 $L^1(Q_T)$。于是在\eqref{align:ch08-8.106} 中令 $k\to\infty$ 并用\eqref{align:ch08-8.108a}，得
\begin{align*}
\int_{\mathbb{R}}f(y)\,\mathrm{d}\nu_{(x,t)}(y)
=\sum_{l=0}^n c_l\,f(s_l(w))+\sum_{m=1}^n d_m\,f(t_m(w))\quad\text{在 }Q_T\text{ 中几乎处处},
\end{align*}
由 $f$ 的任意性即得等式\eqref{align:ch08-8.97}。

只证\eqref{align:ch08-8.108a}，因\eqref{align:ch08-8.108b} 的证明类似。对任意 $\zeta\in L^1(Q_T)$（$T\in(0,\infty)$），有
\begin{align}
\begin{aligned}
&\iint_{Q_T}[\chi_{Q_{k,l}}f((s_l\circ\phi)(v_{\epsilon_k}))-c_l f(s_l(w))]\zeta\,\mathrm{d}x\,\mathrm{d}t\\
&\qquad=\iint_{Q_T}(\chi_{Q_{k,l}}-c_l)f(s_l(w))\zeta\,\mathrm{d}x\,\mathrm{d}t
+\iint_{Q_T}\chi_{Q_{k,l}}[f((s_l\circ\phi)(v_{\epsilon_k}))-f(s_l(w))]\zeta\,\mathrm{d}x\,\mathrm{d}t.
\end{aligned} \label{align:ch08-8.109}
\end{align}
由于 $w\in L^\infty(Q_T)$ 且 $f\circ s_l$ 连续，$f(s_l(w))\in L^\infty(Q_T)$。由\eqref{align:ch08-8.107} 中第一个收敛，
\begin{align}
\lim_{k\to\infty}\iint_{Q_T}(\chi_{Q_{k,l}}-c_l)f(s_l(w))\zeta\,\mathrm{d}x\,\mathrm{d}t=0. \label{align:ch08-8.110}
\end{align}
至于\eqref{align:ch08-8.109} 右端第二个积分，注意由等式\eqref{align:ch08-8.101} 与\cref{proposition:ch05-几乎处处收敛} 有 $\phi(v_{\epsilon_k})\to w$，因而（至多差一个子列）$f((s_l\circ\phi)(v_{\epsilon_k}))\to f(s_l(w))$ 在 $Q_\infty$ 中几乎处处。由控制收敛定理，
\begin{align}
\lim_{k\to\infty}\iint_{Q_T}\chi_{Q_{k,l}}[f((s_l\circ\phi)(v_{\epsilon_k}))-f(s_l(w))]\zeta\,\mathrm{d}x\,\mathrm{d}t=0. \label{align:ch08-8.111}
\end{align}
由\eqref{align:ch08-8.109}--\eqref{align:ch08-8.111} 得\eqref{align:ch08-8.108a} 中的收敛。证毕。
\end{proof}

还需证\cref{proposition:ch08-σ为Dirac}。由\cref{remark:ch08-σ第一等式} 只需证\eqref{align:ch08-8.105} 中第二个等式。我们概述证明，细节见文献~\cite{ref30}。

记 $e_1\leqslant\cdots\leqslant e_{2n}$ 为 $\phi$ 图像的局部极值 $\phi(a_l),\phi(b_l)$（$l=1,\cdots ,n$）。令 $E\equiv E_k\triangleq [e_k,e_{k+1}]$（$k=1,\cdots ,2n-1$）。于是 $e_1=A$，$e_{2n}=B$，且
\begin{align*}
\mathbb{R}=(-\infty,A]\cup\left(\bigcup_{k=1}^{2n-1}E_k\right)\cup[B,\infty).
\end{align*}
对每个 $l=0,\cdots ,n$，$E\cap J_l=\varnothing $ 或 $E\cap J_l=E$；对 $E\cap\widehat J_m$（$m=1,\cdots ,n$）类似。故存在 $p+1$ 个区间 $J_{l_i}$ 与 $p$ 个区间 $\widehat J_{m_j}$（$p=1,\cdots ,n$；$l_i\in\{0,\cdots ,p\}$；$m_j\in\{1,\cdots ,p\}$），使
\begin{align}
E\subseteq\left(\bigcap_{i=0}^p J_{l_i}\right)\cap\left(\bigcap_{j=1}^p\widehat J_{m_j}\right). \label{align:ch08-8.112}
\end{align}
不失一般性设 $l_0\leqslant\cdots\leqslant l_p$、$m_1\leqslant\cdots\leqslant m_p$。借助\cref{proposition:ch08-分离等式}，可证如下引理。

\begin{lemma}[各区间上的矩等式]\label{lemma:ch08-矩等式}
设\cref{theorem:ch08-测度刻画} 的假设满足。则对几乎所有 $(x,t)\in Q_\infty$：
\begin{enumerate}[(i)]
\item\label{item:ch08-mom-i} 设 $K\subseteq(-\infty,A]$（$A\triangleq \min\{\phi(a_1),\cdots ,\phi(a_p)\}$）紧，$\sigma_{(x,t)}(K)>0$。则对几乎所有 $y\in(-\infty,A]$，
\begin{align}
\sigma^0_{(x,t)}((-\infty,y])\,\sigma_{(x,t)}(K)=\sigma^0_{(x,t)}((-\infty,y]\cap K); \label{align:ch08-8.113}
\end{align}
\item\label{item:ch08-mom-ii} 设 $K\subseteq[B,\infty)$（$B\triangleq \max\{\phi(b_1),\cdots ,\phi(b_p)\}$）紧，$\sigma_{(x,t)}(K)>0$。则对几乎所有 $y\in[B,\infty)$，
\begin{align}
\sigma^n_{(x,t)}([y,\infty))\,\sigma_{(x,t)}(K)=\sigma^n_{(x,t)}([y,\infty)\cap K); \label{align:ch08-8.114}
\end{align}
\item\label{item:ch08-mom-iii} 设 $E\equiv[e_k,e_{k+1}]$（$k=1,\cdots ,2n-1$），$K\subseteq E$ 紧，$\sigma_{(x,t)}(K)>0$。则对几乎所有 $y\in E$，
\begin{align*}
\sum_{i=1}^{p+1}D_i(y)\left\{\frac{\sigma^{l_i}_{(x,t)}(K)}{\sigma_{(x,t)}(K)}-\sigma^{l_i}_{(x,t)}(E)\right\}
+\sum_{j=1}^p D_{j+1}(y)\left\{\frac{\widehat\sigma^{m_j}_{(x,t)}(K)}{\sigma_{(x,t)}(K)}-\widehat\sigma^{m_j}_{(x,t)}(E)\right\}\qquad\nonumber\\
=\sum_{i=1}^{p+1}s'_{l_i}(y)\left\{\rho^{l_i}_{(x,t)}(y)-\frac{\rho^{l_i}_K(y)}{\sigma_{(x,t)}(K)}\right\}
+\sum_{j=1}^p t'_{m_j}(y)\left\{\widehat\rho^{m_j}_{(x,t)}(y)-\frac{\widehat\rho^{m_j}_K(y)}{\sigma_{(x,t)}(K)}\right\},
\end{align*}
其中
\begin{align*}
D_1(y)\triangleq 0,\quad D_i(y)\triangleq \sum_{k=1}^{i-1}(s'_{l_k}-t'_{m_k})(y)\quad(i=2,\cdots ,p+1),
\end{align*}
\begin{align*}
\rho^{l_i}_{(x,t)}(y)\triangleq \sigma^{l_i}_{(x,t)}([y,\phi(b_{l_i+1})]),\quad
\widehat\rho^{m_j}_{(x,t)}(y)\triangleq \widehat\sigma^{m_j}_{(x,t)}([y,\phi(b_{m_j})]),
\end{align*}
\begin{align*}
\rho^{l_i}_K(y)\triangleq \sigma^{l_i}_{(x,t)}([y,\phi(b_{l_i+1})]\cap K),\quad
\widehat\rho^{m_j}_K(y)\triangleq \widehat\sigma^{m_j}_{(x,t)}([y,\phi(b_{m_j})]\cap K).
\end{align*}
\end{enumerate}
\end{lemma}

\begin{proof}
只证 (i)--(ii)，(iii) 的证明类似但更繁（见文献~\cite{ref30} 引理 5.5）。由于 $K\subseteq(-\infty,A]$ 紧，存在序列 $\{f_h\}\subseteq C_c(\mathbb{R})$ 使 $0\leqslant f_h\leqslant1$、$f_h=1$ 在 $K$ 上，且 $\lim_{h\to\infty}f_h(y)=\chi_K(y)$ 对所有 $y\in\mathbb{R}$。令 $a_0\in I_0$ 为使 $\phi(a_0)=A$ 的唯一点，取\eqref{align:ch08-8.14} 中的 $k$ 使 $G_\phi(y)=\int_{a_0}^y g(\phi(z))\,\mathrm{d}z$（$y\in\mathbb{R}$），其中 $g\in C^1_c(\mathbb{R})$ 且 $\operatorname{supp}g\subseteq(-\infty,A]$。由等式\eqref{align:ch08-8.58} 有
\begin{align*}
\int_{\mathbb{R}}f_h(\phi(y))G_\phi(y)\,\mathrm{d}\nu_{(x,t)}(y)
=\int_{\mathbb{R}}f_h(\phi(y))\,\mathrm{d}\nu_{(x,t)}(y)\int_{\mathbb{R}}G_\phi(y)\,\mathrm{d}\nu_{(x,t)}(y),
\end{align*}
由 $g$ 的选取与\cref{remark:ch08-σ第一等式} 得
\begin{align*}
\int_{(-\infty,A]}f_h(y)(G_\phi\circ s_0)(y)\,\mathrm{d}\sigma^0_{(x,t)}(y)
=\int_{\phi(\mathbb{R})}f_h(y)\,\mathrm{d}\sigma_{(x,t)}(y)\int_{(-\infty,A]}(G_\phi\circ s_0)(y)\,\mathrm{d}\sigma^0_{(x,t)}(y).
\end{align*}
令 $h\to\infty$，得
\begin{align*}
\int_K(G_\phi\circ s_0)(y)\,\mathrm{d}\sigma^0_{(x,t)}(y)
=\sigma_{(x,t)}(K)\int_{(-\infty,A]}(G_\phi\circ s_0)(y)\,\mathrm{d}\sigma^0_{(x,t)}(y),
\end{align*}
因而显然
\begin{align*}
\int_{(-\infty,A]}\mathrm{d}\sigma^0_{(x,t)}(y)\int_A^y g(z)s_0'(z)\,\mathrm{d}z
=\frac{1}{\sigma_{(x,t)}(K)}\int_K\mathrm{d}\sigma^0_{(x,t)}(y)\int_A^y g(z)s_0'(z)\,\mathrm{d}z.
\end{align*}
交换积分顺序即得\eqref{align:ch08-8.113}。情形 (ii) 类似地得
\begin{align*}
\int_{[B,\infty)}\mathrm{d}\sigma^n_{(x,t)}(y)\int_B^y g(z)s_0'(z)\,\mathrm{d}z
=\frac{1}{\sigma_{(x,t)}(K)}\int_K\mathrm{d}\sigma^n_{(x,t)}(y)\int_B^y g(z)s_0'(z)\,\mathrm{d}z,
\end{align*}
由上式得\eqref{align:ch08-8.114}。证毕。
\end{proof}

现可证\cref{proposition:ch08-σ为Dirac}。

\begin{proof}[\cref{proposition:ch08-σ为Dirac} 的证明]
设 $\operatorname{supp}\sigma_{(x,t)}|_{(-\infty,A]}\neq\varnothing $，令
\begin{align*}
q_0\equiv q_0(x,t)\triangleq \max\{y\in\operatorname{supp}\sigma_{(x,t)}|_{(-\infty,A]}\}.
\end{align*}
对任意 $\delta>0$ 考察紧集 $K=[q_0-\delta,q_0]$。则 $\sigma_{(x,t)}(K)>0$，且对所有 $y\in(-\infty,q_0-\delta]$ 有
\begin{align*}
\sigma^0_{(x,t)}((-\infty,y]\cap K)=0.
\end{align*}
由等式\eqref{align:ch08-8.113} 与 $\delta$ 的任意性得 $\sigma^0_{(x,t)}((-\infty,q_0))$，因而 $\operatorname{supp}\sigma_{(x,t)}|_{(-\infty,A]}=\{q_0\}$。进而取 $K=\{q_0\}$ 时等式\eqref{align:ch08-8.113} 给出 $\sigma^0_{(x,t)}(\{q_0\})=\sigma_{(x,t)}(\{q_0\})=1$。用\eqref{align:ch08-8.114} 类似地可见 $\operatorname{supp}\sigma_{(x,t)}|_{[B,\infty)}$ 至多由一点组成，记为 $\{q_{2n}\}$；若非空，则 $\sigma^n_{(x,t)}(\{q_{2n}\})=\sigma_{(x,t)}(\{q_{2n}\})=1$。借助\cref{lemma:ch08-矩等式}(iii) 的更精细论证表明，对每个区间 $E_k=[e_k,e_{k+1}]$（$k=1,\cdots ,2n-1$）同样成立：若 $\operatorname{supp}\sigma|_{E_k}\neq\varnothing $，则 $\operatorname{supp}\sigma_{(x,t)}|_{E_k}=\{q_k\}$，且对某个 $q_k\in E_k$ 有 $\sigma_{(x,t)}(\{q_k\})=1$（细节见文献~\cite{ref30} 命题 5.2 的证明）。

综上，测度 $\sigma_{(x,t)}$ 的支撑至多由 $2n+1$ 个点 $\{q_k\}$ 组成，且每个 $k=0,\cdots ,2n$ 有 $\sigma_{(x,t)}(\{q_k\})=1$。然而由于 $\sigma_{(x,t)}\in\mathcal{P}(\mathbb{R})$，这蕴含 $\operatorname{supp}\sigma_{(x,t)}$ 由一点组成，记为 $q=q(x,t)$，且 $\sigma_{(x,t)}=\delta_{q(x,t)}$。由\eqref{align:ch08-8.105} 中第一个等式有
\begin{align*}
q(x,t)=\int_{\mathbb{R}}y\,\mathrm{d}\sigma_{(x,t)}(y)
=\int_{\mathbb{R}}\phi(y)\,\mathrm{d}\nu_{(x,t)}(y)=\phi^*(x,t),
\end{align*}
由\cref{remark:ch08-w等于phi星} 得 $\sigma_{(x,t)}=\delta_{\phi^*(x,t)}=\delta_{w(x,t)}$ 对几乎所有 $(x,t)\in Q_\infty$。证毕。
\end{proof}




\end{document}